\section{Polynomial Commitment Scheme}
\label{sec:pcs}

We give two instantiations of the commitment interface of
Section~\ref{sec:highlevelPCS}. Both bind a matrix
\(\bm A\in\Field^{m\times\kappa}\), where \(N=m\kappa\) and
\([\ExtensionField:\Field]=\kappa\).
The first starts from a multilinear PCS over \(\ExtensionField\)
and reduces a packed query to one terminal MLE opening.
The second starts from a linear code over \(\Field\): scalar
extension gives a code for the packed vector, and code switching
to a BaseFold code supplies the extension-field PCS needed by
the first construction.

\subsection{Given a PCS over Extension Field}
\label{sec:pcs-extension}

\medskip\noindent{\bf Commitment and Opening Relation}

Let \(m,\kappa\) be powers of two, with \(\kappa\geq2\).
Fix an \(\Field\)-basis
\(\ext{\bm\beta}=(\ext\beta_1,\ldots,\ext\beta_\kappa)^\top\)
of \(\ExtensionField\), and let \(\pc\) be a multilinear PCS for
vectors in \(\ExtensionField^m\).
We represent extension-field elements in this basis, so coordinate
extraction and basis packing are operations on their representation.
The schoolbook specialization below uses a polynomial basis for this
representation.

The basis and all coordinate matrices are defined over the same
working field as the repeated query-side checks.
Section~\ref{sec:binary-packing} uses this interface for the
binary-field instantiation.

\paragraph{Setup and commitment.}
Run \(\pc.\Setup(1^\lambda,m)\) to obtain
\(\mathsf{pp}_{\pc}\).
The public parameters also contain the fields, the basis, \(N,m,\kappa\),
and the repetition count \(k\) specified below.
Arrange the witness as a matrix and pack each row into one
extension-field element:
\[
    \bm A=\concat{\bm a^{(1)},\ldots,\bm a^{(\kappa)}}
        \in\Field^{m\times\kappa},
    \qquad
    \ext{\bm a}=\bm A\ext{\bm\beta}\in\ExtensionField^m.
\]
The commitment is
\(\mathsf{cm}=\pc.\Com(\mathsf{pp}_{\pc},\ext{\bm a};
\rho_{\pc})\), where \(\rho_{\pc}\) is the underlying commitment
randomness. The prover retains \(\bm A\) and \(\rho_{\pc}\).
The fixed basis makes the packing injective on
\(\Field^{m\times\kappa}\).

\paragraph{Opening relation.}
Following the overview, the public query consists of packing weights
\(\ext{\bm\rho}\in\ExtensionField^\kappa\) and a tensor vector
\(\ext{\bm\chi}\in\ExtensionField^m\). The claimed value is
\(\ext u\in\ExtensionField\):
\begin{equation}
    \ext u
    =\langle\bm A\ext{\bm\rho},\ext{\bm\chi}\rangle
    =\ext{\bm\rho}^{\top}\bm A^\top\ext{\bm\chi}.
\label{eq:pcs-opening-relation}
\end{equation}
The query is supplied by its two-entry tensor factors:
\[
    \ext{\bm\chi}
    =\bigotimes_{j=1}^{\log_2m}\ext{\bm\chi}^{(j)},
    \qquad
    \ext{\bm\chi}^{(j)}\in\ExtensionField^2.
\]
The factors need not be normalized MLE weights. This includes the
weighted Hadamard query: coordinatewise products of tensor vectors
are obtained by multiplying their corresponding factors coordinatewise.
For an ordinary MLE opening of the length-\(N\) witness, both
\(\ext{\bm\rho}\) and \(\ext{\bm\chi}\) are tensor spans of the
corresponding parts of the evaluation point.

\medskip\noindent{\bf Packed Opening Protocol}

\paragraph{Reduce the column claims over \(\ExtensionField\).}
Write the public query and the claimed column values in base-field
coordinates:
\[
    \ext{\bm\chi}=\bm\Xi\ext{\bm\beta},
    \qquad
    \ext{\bm v}=\bm V\ext{\bm\beta},
    \qquad
    \bm\Xi\in\Field^{m\times\kappa},
    \quad
    \bm V\in\Field^{\kappa\times\kappa}.
\]
The honest prover computes \(\bm V=\bm A^\top\bm\Xi\) and sends
\(\ext{\bm v}=\bm V\ext{\bm\beta}\).
The verifier recovers \(\bm V\) from this message and checks
\(\ext u=\ext{\bm\rho}^{\top}\ext{\bm v}\).
It does not construct \(\bm\Xi\).

Since \(\bm A\) has base-field entries, the remaining claim is
equivalent to a base-field matrix identity:
\begin{equation*}
    \ext{\bm v}=\bm A^\top\ext{\bm\chi}
    \quad\Longleftrightarrow\quad
    \bm V=\bm A^\top\bm\Xi.
\end{equation*}
Packing the columns of this identity is injective and exposes the
committed vector:
\begin{equation}
    \bm V=\bm A^\top\bm\Xi
    \quad\Longleftrightarrow\quad
    \ext{\bm\beta}^{\top}\bm V
    =\ext{\bm a}^{\top}\bm\Xi.
\label{eq:pcs-encoded-matrix}
\end{equation}
After \(\ext{\bm v}\), and hence \(\bm V\), is fixed, the verifier
samples \(\ext{\bm r}\sample\ExtensionField^\kappa\) uniformly.
It computes \(\ext c=\ext{\bm\beta}^{\top}\bm V\ext{\bm r}\)
and reduces Eq.~\eqref{eq:pcs-encoded-matrix} to:
\begin{equation}
    \ext c
    =\ext{\bm\beta}^{\top}\bm V\ext{\bm r}
    =\langle\ext{\bm a},\bm\Xi\ext{\bm r}\rangle.
\label{eq:pcs-random-opening}
\end{equation}
Here \(\ext{\bm\rho}\) belongs to the original query, whereas
\(\ext{\bm r}\) is a fresh projection challenge.

Apply Fact~\ref{fact:ip} over \(\ExtensionField\) to
Eq.~\eqref{eq:pcs-random-opening}.
Let \(\ext{\bm y}\in\ExtensionField^{\log_2m}\) be its terminal
point and set
\(\ext{\bm\xi}=\tensorspan{\ext{\bm y}}\).
The prover's terminal values
\(\ext w_{\mathsf{com}},\ext w\in\ExtensionField\) give two claims:
\begin{equation}
    \ext w_{\mathsf{com}}
    =\langle\ext{\bm a},\ext{\bm\xi}\rangle,
    \qquad
    \ext w=\ext{\bm\xi}^{\top}\bm\Xi\ext{\bm r}.
\label{eq:pcs-terminal-claims}
\end{equation}
The first is an MLE opening of the committed vector.
We verify the second using the public tensor queries, without
materializing their coordinate matrices at the verifier.

\paragraph{Reduce the induced query over \(\Field\).}
Decompose the terminal tensor weight as
\(\ext{\bm\xi}=\bm U\ext{\bm\beta}\), with
\(\bm U\in\Field^{m\times\kappa}\).
The prover sends a claimed matrix product
\(\bm W\in\Field^{\kappa\times\kappa}\).
The required identities are:
\begin{equation}
    \bm W=\bm U^\top\bm\Xi,
    \qquad
    \ext w=\ext{\bm\beta}^{\top}\bm W\ext{\bm r}.
\label{eq:pcs-query-matrix}
\end{equation}
The verifier checks the scalar identity directly.
It verifies the matrix identity by \(k\) independent repetitions.
In repetition \(i\), after \(\bm W\) is fixed, the verifier samples
independent uniform
\(\bm s_{(i)},\bm t_{(i)}\in\Field^\kappa\).
The parties apply Fact~\ref{fact:ip}, now over \(\Field\), to:
\begin{equation}
    \bm s_{(i)}^\top\bm W\bm t_{(i)}
    =\langle\bm U\bm s_{(i)},\bm\Xi\bm t_{(i)}\rangle.
\label{eq:pcs-base-projection}
\end{equation}
All challenges in this inner-product reduction lie in \(\Field\).
Its terminal point
\(\bm y_{(i)}\in\Field^{\log_2m}\) determines
\(\bm l_{(i)}=\tensorspan{\bm y_{(i)}}\in\Field^m\).
Let \(v_{(i),1},v_{(i),2}\in\Field\) be its terminal values.

The verifier computes the following two public tensor contractions:
\begin{equation}
\begin{aligned}
    \langle\bm l_{(i)},\ext{\bm\xi}\rangle
       &=(\bm l_{(i)}^\top\bm U)\ext{\bm\beta},\\
    \langle\bm l_{(i)},\ext{\bm\chi}\rangle
       &=(\bm l_{(i)}^\top\bm\Xi)\ext{\bm\beta}.
\end{aligned}
\label{eq:pcs-public-query-evaluation}
\end{equation}
Since \(\bm l_{(i)}\) lies in \(\Field^m\), the two rows on the
right lie in \(\Field^{1\times\kappa}\). They are therefore exactly
the basis coordinates of the two computed values.
Using these recovered rows, the verifier checks:
\begin{equation}
    v_{(i),1}=(\bm l_{(i)}^\top\bm U)\bm s_{(i)},
    \qquad
    v_{(i),2}=(\bm l_{(i)}^\top\bm\Xi)\bm t_{(i)}.
\label{eq:pcs-implicit-query-check}
\end{equation}
No commitment to \(\bm U\) or \(\bm\Xi\) is needed: both are
determined by public queries, and these terminal checks are exact.

The tensor factors make Eq.~\eqref{eq:pcs-public-query-evaluation}
efficient to evaluate. For example:
\[
    \langle\bm l_{(i)},\ext{\bm\chi}\rangle
    =
    \prod_{j=1}^{\log_2m}
    \left\langle
        (1-y_{(i),j},y_{(i),j})^\top,
        \ext{\bm\chi}^{(j)}
    \right\rangle.
\]
The same factorization applies to \(\ext{\bm\xi}\).
Each contraction uses \(O(\log m)\) mixed-field and extension-field
operations, rather than a length-\(m\) vector evaluation.

\paragraph{Repetition and terminal opening.}
Assume \(|\Field|>2\log_2m+2\), and set:
\begin{equation}
    k=
    \left\lceil
        \frac{\lambda}{
            \log_2\!\left(
                |\Field|/(2\log_2m+2)
            \right)
        }
    \right\rceil.
\label{eq:pcs-repetition-number}
\end{equation}
Each repetition uses fresh projection and inner-product challenges,
while the same matrix \(\bm W\) remains fixed.
Finally, the parties invoke the underlying opening protocol for
\(\ext w_{\mathsf{com}}\) at \(\ext{\bm y}\), using the interface
of Definition~\ref{def:pcs}:
\[
    \pc.\Open\!\left(
        \mathsf{pp}_{\pc},\mathsf{cm},
        \ext w_{\mathsf{com}},\ext{\bm y}
        :\ext{\bm a},\rho_{\pc}
    \right).
\]
The verifier accepts only if every scalar check, every inner-product
reduction, all checks in Eq.~\eqref{eq:pcs-implicit-query-check},
and this terminal opening accept.
Algorithm~\ref{alg:packed-pcs} in Appendix~\ref{appx:pseudocode}
summarizes the opening protocol.

\medskip\noindent{\bf Security and Complexity}

\begin{proposition}[Packed polynomial commitment scheme]
\label{proposition:packed-pcs}
Suppose that \(\pc\) binds the prover to a vector in
\(\ExtensionField^m\) and has evaluation-soundness error
\(\epsilon_{\mathsf{PCS}}\) for adaptively chosen terminal MLE
claims relative to that vector.
Under the parameters above, the packed commitment is binding and
the opening protocol is complete.
Its evaluation-soundness error is at most:
\begin{equation}
    \epsilon_{\mathsf{PCS}}
    +\frac{1+2\log_2m}{|\ExtensionField|}
    +\left(\frac{2\log_2m+2}{|\Field|}\right)^k.
\label{eq:pcs-soundness}
\end{equation}
With \(k\) from Eq.~\eqref{eq:pcs-repetition-number}, the last term
is at most \(2^{-\lambda}\).

Excluding the underlying PCS opening, the prover's arithmetic cost is:
\begin{equation}
\begin{aligned}
    &O\bigl((\kappa+k)N+k\kappa^2\bigr)\,\bb\\
    &\qquad{}+O(N)\,\be+O(m+\kappa)\,\ee.
\end{aligned}
\label{eq:pcs-prover-cost}
\end{equation}
The verifier's arithmetic cost is:
\begin{equation}
\begin{aligned}
    &O\bigl(k(\kappa^2+\log m)\bigr)\,\bb
       +O(k\log m)\,\be\\
    &\qquad{}+O\bigl(\kappa+(k+1)\log m\bigr)\,\ee.
\end{aligned}
\label{eq:pcs-verifier-cost}
\end{equation}
When \(m\geq\kappa\), schoolbook arithmetic reduces
Eq.~\eqref{eq:pcs-prover-cost} to
\(O((\kappa+k)N)\,\bb\).
The opening reduction uses \(O(N+\kappa^2)\) base-field coordinates
of peak prover memory, or \(O(N)\) when \(m\geq\kappa\), excluding
the workspace of the underlying PCS.

With \(\ext u\) already part of the public statement, prover-to-verifier
communication, excluding the underlying opening, consists of:
\[
\begin{aligned}
    &\kappa+2\log_2m+2
        &&\text{elements of }\ExtensionField,\\
    &\kappa^2+k(2\log_2m+2)
        &&\text{elements of }\Field.
\end{aligned}
\]
If the underlying commitment is hiding, the packed commitment is also
hiding. This statement concerns the commitment, not the unmasked
opening transcript, which reveals the column evaluations
\(\ext{\bm v}\).
\end{proposition}

The proof is deferred to Appendix~\ref{appx:proof-pcs}.

\subsection{Given a linear code on Base Field}
\label{sec:pcs-base-code}

We use the code-switching construction of
Blaze~\cite{eurocrypt/BrehmCFRRZ25}, with a BaseFold
code~\cite{crypto/ZeilbergerCF24} as the target, under its
source-encoding and proximity assumptions.

Let \(\Code:\Field^m\to\Field^n\) be a linear code with
generator matrix \(\bm G\in\Field^{n\times m}\).
Its \(\kappa\)-way interleaving encodes the columns of
\(\bm A\in\Field^{m\times\kappa}\) independently.
Using the fixed basis from Section~\ref{sec:pcs-extension},
the resulting codeword has the equivalent packed representation:
\begin{equation}
\begin{aligned}
    \bm C&=\bm G\bm A\in\Field^{n\times\kappa},\\
    \ext{\bm c}&=\bm C\ext{\bm\beta}
       =\bm G(\bm A\ext{\bm\beta})
       =\Code_{\ExtensionField}(\ext{\bm a}).
\end{aligned}
\label{eq:pcs-code-packing}
\end{equation}
Here \(\Code_{\ExtensionField}\) is the scalar extension defined
by the same generator matrix; it has the same rate and relative
distance as \(\Code\).
Thus each row of the interleaved codeword is one symbol of an
extension-field codeword for \(\ext{\bm a}=\bm A\ext{\bm\beta}\).
Committing to the rows of \(\bm C\) supplies the source codeword
oracle required by Blaze.

Apply Blaze's code-switching construction over \(\ExtensionField\)
to this oracle, switching to a BaseFold encoding of the same
length-\(m\) message \(\ext{\bm a}\).
The resulting MLE openings instantiate the underlying PCS \(\pc\)
in Section~\ref{sec:pcs-extension}.
That subsection handles the challenge-dependent packing weights
\(\ext{\bm\rho}\); code switching authenticates its terminal
fixed-basis opening.
The switching costs and soundness error are those of the chosen
Blaze instantiation, including its source-encoding consistency
argument, and are included in the underlying PCS costs and
\(\epsilon_{\mathsf{PCS}}\), respectively.

If \(\Code\) is constant-rate and linear-time encodable, forming
\(\bm C\) costs \(O(N)\) base-field operations and its payload
occupies \(O(N\log_2|\Field|)\) bits, excluding Merkle-tree hashes.
The BaseFold stage works on \(m=N/\kappa\) extension-field
message symbols. This gives the bit-width-aware instantiation:
encode the interleaved witness over the small field, then switch
the packed codeword to the opening-friendly code over the extension.
